
\section{A second-order Leray correction and a meridional dual witness}
\label{sec:second-order-leray}
This section concerns a smooth background and does not establish unforced
finite-time blowup. Its short-time scheme is an explicit instance of standard
Leray--Stokes perturbation theory, not an originality claim.
Let $s>5/2$, $\nu>0$, and let $u_B$ be a smooth divergence-free field on
$\mathbb T^3$, with complete residual
\[
R_B=\partial_tu_B-\nu\Delta u_B+(u_B\cdot\nabla)u_B+\nabla p_B.
\]
Let $\mathbb P$ denote the Leray projection,
$S(t)=e^{\nu t\Delta}$ and
$B(a,b)=\mathbb P[(a\cdot\nabla)b]$.
For a slab $[t_0,t_0+T]$ put
$A=\sup_t\|u_B(t)\|_{H^s}$,
$F=\sup_t\|\mathbb P R_B(t)\|_{H^s}$.
Define
\begin{align*}
w_1(t)&=-\int_{t_0}^tS(t-v)\mathbb P R_B(v)\,dv,\\
N_1&=B(u_B,w_1)+B(w_1,u_B)+B(w_1,w_1),\\
w_2(t)&=-\int_{t_0}^t S(t-v)N_1(v)\,dv.
\end{align*}
\begin{proposition}[Exact residual after two projected corrections]
With $w=w_1+w_2$, the projected unforced residual is precisely
\[
\mathcal E_2=B(u_B,w_2)+B(w_2,u_B)+B(w_1,w_2)
+B(w_2,w_1)+B(w_2,w_2).
\]
If $\|B(a,b)\|_{H^{s-1}}\le C_s\|a\|_{H^s}\|b\|_{H^s}$,
and
\[
H_\nu(T)=T+\sqrt{\frac{2T}{e\nu}},\quad
W_1=TF,\quad W_2=C_sH_\nu(T)(2AW_1+W_1^2),
\]
then
\[
\|w_1\|_{CH^s}\le W_1,\quad
\|w_2\|_{CH^s}\le W_2,\quad
\|\mathcal E_2\|_{CH^{s-1}}
\le C_s(2AW_2+2W_1W_2+W_2^2).
\]
The full unprojected residual equals $\mathcal E_2$ upon taking
\[
\pi=-\Delta^{-1}\nabla\cdot
\{R_B+(u_B\cdot\nabla)w+(w\cdot\nabla)u_B+(w\cdot\nabla)w\}
\]
as the correction pressure, with zero spatial mean.
\end{proposition}
\begin{proof}
The defining Duhamel equations give
$(\partial_t-\nu\Delta)w_1=-\mathbb P R_B$ and
$(\partial_t-\nu\Delta)w_2=-N_1$. Expand the full quadratic
nonlinearity and cancel the terms containing $w_1$ alone.
The Fourier multiplier inequality
$\sqrt{1+|k|^2}e^{-\nu t|k|^2}
\le1+(2e\nu t)^{-1/2}$
and time integration prove the stated $H_\nu(T)$ bound.
The bilinear estimate then yields all three inequalities.
Finally, the Poisson pressure removes the gradient component of the
total forcing; crucially, it retains $\nabla\cdot R_B$.
\end{proof}
For fixed $A,F$ the resulting residual bound is $O(T^{3/2})$;
there is no claim of exact zero residual or terminal continuation.

\begin{proposition}[Explicit meridional dual witness on $\mathbb R^3$]
Let $u_B=r e^{-(r^2+z^2)}e_\theta$ be a stationary Gaussian
pure-swirl background. Define
$\chi(x)=(x-\frac12)^4(1-x)^4$ for $x\in[\frac12,1]$ and
$\chi(x)=0$ otherwise, and put
$\psi(r,z)=\chi(r)\chi(z)$, $\Phi=\nabla\times(\psi e_\theta)$.
Then $\Phi$ is compactly supported and divergence-free, and
\[
\langle R_B,\Phi\rangle
=8\pi\int_{1/2}^1\!\int_{1/2}^1
r^2 z e^{-2(r^2+z^2)}\chi(r)\chi(z)\,dr\,dz
>\frac{\pi e^{-4}}{322560^2}>0.
\]
Consequently a smooth unforced correction vanishing at $t_0$
must satisfy
$\langle\partial_t w(t_0),\Phi\rangle
<-\pi e^{-4}/322560^2$.
\end{proposition}
\begin{proof}
The meridional centrifugal residual has
$(R_B)_r=-r e^{-2(r^2+z^2)}$, $(R_B)_z=0$
modulo any pressure gradient. Thus
$(\nabla\times R_B)_\theta=4rz e^{-2(r^2+z^2)}$.
Use integration by parts, the cylindrical volume element
$2\pi r\,dr\,dz$, and the exact beta integral
$\int_{1/2}^1\chi(x)\,dx=1/322560$.
The initial-time identity follows by testing the correction equation
with the solenoidal $\Phi$.
\end{proof}
The Gaussian witness is a pressure-independent \emph{positive control},
not a computation on the source's fully corrected field.
